\section{Automatic Commit-Reveal Insertion}
\label{sec:commit-reveal}

This section describes the EventGraph rewrite that protects
simultaneous public moves from transaction-ordering adversaries.
The pass is small but the property it secures is one of the most
visible payoffs of using Vegas: a programmer who writes
``two players \TT{yield} simultaneously'' gets the commit-reveal
protocol the textbook would have insisted on, deployed on chain
with no further work.

\subsection{The problem}
\label{sec:cr:problem}

Consider the two-line specification
\begin{lstlisting}[style=vg]
yield A(x: bool);
yield B(y: bool);
\end{lstlisting}
At the EventGraph level these are two PUBLIC writes on concurrent
nodes.  Read as a game, the analysis is unambiguous: A and B move
simultaneously without observing each other.  Read as a Solidity
contract, the implementation has to be cautious.  If A's
transaction is in the public mempool, B can read \TT{x} before
deciding \TT{y}.  An ordering adversary (a builder, or a competing
searcher) can promote the read-then-write transaction.  The
specification's notion of simultaneity is broken by the deployment
target's notion of transaction ordering.

This is a specific subclass of MEV --- \emph{observation-before-action}
on simultaneous moves --- and the canonical mitigation is well
known~\citep{Blum1983CoinFlipping,Pedersen1991Commitment,breidenbach2018enter}:
each role first commits to a hash of their move, and only after
all roles have committed does anybody reveal.  Vegas applies the
mitigation automatically, by rewriting the EventGraph.

\subsection{Risk partners}

Two nodes are \emph{risk partners} when both are pure-public
writes, neither reaches the other in the dependency DAG, and they
are distinct.  This is precisely the configuration in which the
ordering adversary problem arises.  The pass identifies the set
of nodes that have at least one risk partner --- call this set
$\mathrm{Split}$ --- and rewrites every node in $\mathrm{Split}$.

Pure-public writes (every visibility tag is \PUBLIC) are the only
candidates for rewriting.  Nodes that already commit (visibility
\COMMIT) need no rewriting; their reveal partner is already
present.  Nodes that themselves have a guard reading another
actor's private value cannot arise (the typechecker rejects such
reads).  And nodes that are in a strict dependency relation with
their potential partner are not concurrent, so the ordering
adversary has no informational asymmetry to exploit.

\subsection{The rewrite}
\label{sec:cr:rewrite}

For each $a \in \mathrm{Split}$ the pass introduces two fresh
nodes:

\begin{itemize}\itemsep1pt
\item a \emph{commit node} $a^{c}$, replacing $a$ as the new owner
  of the deposit (if $a$ was a \TT{join} step) and of the
  original predecessors of $a$, but with all of $a$'s \PUBLIC
  writes retagged as \COMMIT and with a trivial guard
  (anything is allowed to commit);
\item a \emph{reveal node} $a^{r}$, carrying $a$'s original guard,
  with all the same writes retagged as \REVEAL.
\end{itemize}

Edges are wired so that

\begin{itemize}\itemsep1pt
\item $a^{c}$ has the same predecessors $a$ originally had;
\item $a^{r}$ has predecessors $\{a^{c}\} \cup \{b^{c} : b \in
  \mathrm{partners}(a)\} \cup \mathrm{pred}(a)$;
\item every node that previously depended on $a$ now depends on
  $a^{r}$ instead.
\end{itemize}

The second clause is the informational barrier: $a^{r}$ cannot
fire until \emph{all} of $a$'s risk partners have also committed.
After the rewrite, no \REVEAL node is on the frontier of any
configuration in which a sibling \COMMIT is still unresolved.
\Cref{sec:eventgraph:frontier} called this the structural
information barrier; the rewrite is what guarantees it for any
group of originally-simultaneous public moves.

\subsection{Illustration}

The bidder fragment of the Vickrey auction
(\S\ref{sec:lang:vickrey}) is the cleanest case.  The source

\begin{lstlisting}[style=vg]
yield or split B1(b: bid) B2(b: bid) B3(b: bid);
\end{lstlisting}

produces, before the rewrite, three concurrent \PUBLIC nodes
(left) and after the rewrite, three commit nodes followed by three
reveal nodes that all wait for the full commit phase (right).

\begin{center}
\begin{tikzpicture}[
  every node/.style={font=\small},
  pu/.style={draw, rounded corners, inner sep=2pt, minimum width=1.5cm, minimum height=4mm, fill=blue!8},
  cnode/.style={draw, inner sep=2pt, minimum width=1.5cm, minimum height=4mm, fill=red!10},
  rv/.style={draw, regular polygon, regular polygon sides=6, inner sep=0pt,
             minimum width=1.6cm, minimum height=5mm, fill=green!10},
  arr/.style={-{Stealth[length=2mm]}, thick}
]
\node[pu] (b1) at (-0.8, 0)   {B1.b};
\node[pu] (b2) at ( 0.6, 0)   {B2.b};
\node[pu] (b3) at ( 2.0, 0)   {B3.b};
\node     at ( 0.6,-0.9) {\textit{before rewrite}};

\node[cnode] (c1) at ( 6.0, 0.5) {B1.b};
\node[cnode] (c2) at ( 7.4, 0.5) {B2.b};
\node[cnode] (c3) at ( 8.8, 0.5) {B3.b};
\node[rv] (r1) at ( 6.0,-1.2) {B1.b};
\node[rv] (r2) at ( 7.4,-1.2) {B2.b};
\node[rv] (r3) at ( 8.8,-1.2) {B3.b};
\node     at ( 7.4,-2.1) {\textit{after rewrite}};
\foreach \c in {c1, c2, c3}
  \foreach \r in {r1, r2, r3}
    \draw[arr, gray] (\c) -- (\r);
\end{tikzpicture}
\end{center}

After the rewrite, each bidder's reveal depends on every bidder's
commit.  Reading off the new graph: no reveal is enabled on a
frontier where any commit is missing, so no bidder learns another
bidder's value before they themselves have committed.

\subsection{What changes for the backends}

Every backend sees the rewritten graph, not the original.  The
operational backends (Solidity/Vyper) therefore emit a
commit-then-reveal pair of contract functions per rewritten
action; the Solidity backend uses the hidden write's payload as
the input to a hash and stores the hash, then on reveal verifies
the hash and assigns the public field.  The Gambit backend reads
the rewritten graph as a sequence of two simultaneous-move rounds
(commits, then reveals); information sets are computed from the
new graph in the usual way, and they correctly reflect the
informational barrier.  No backend-specific commit-reveal logic
exists: the property is paid for once, in the IR, and consumed
everywhere.

\subsection{What the rewrite does \emph{not} do}

A few caveats.

\paragraph{Cryptographic strength is not the IR's concern.}  The
rewrite stipulates \emph{that} a hidden value is committed before
the public reveal; it does not specify the cryptographic primitive.
The Solidity backend chooses a domain-separated keccak with a
random salt and a per-game tag; a different backend (Bitcoin) uses
its native opcodes.  The choice is the backend's; the IR-level
property holds for any binding, hiding commitment.

\paragraph{The rewrite does not address adaptive adversaries
elsewhere.}  An adversary can still censor a player's transaction
to force a timeout; the bail-out subgame (\S\ref{sec:liveness})
makes this a strategic choice with explicit payoff, so the game
model sees it, but censorship-resistance per se is not what the
commit-reveal pass provides.

\paragraph{The rewrite does not enlarge the set of safe
patterns.}  Only \emph{originally-concurrent} public writes are
rewritten.  A public write that is in strict order with all other
public writes is not at risk and stays a single node.  This is
intentional: there is no point in adding round-trips for moves
that the dependency relation already serialises.

\subsection{The information-barrier property}
\label{sec:cr:barrier}

What the rewrite produces is a graph in \emph{commit-reveal
normal form}: every node whose original write was concurrent
with another's has been split into a commit-then-reveal pair,
and every reveal in the result is dominated in the DAG by all
of its risk partners' commits.  The strategic property the
rewrite delivers is the \emph{information barrier}: in any
configuration of the rewritten graph, no reveal node is on the
frontier while any of its commit prerequisites --- its own or a
risk-partner's --- is still unresolved.

Stated informally for the rewritten graph $G$:
\begin{quote}\itshape
For every reveal node $r$ and every configuration $C$ of $G$
with $r$ enabled in $C$, every node $c \in \mathrm{commitPred}(r)$
has been assigned in $C$.
\end{quote}
By the construction of the rewrite this holds: each
$r = a^r$ has $\{a^c\} \cup \{b^c : b \in \mathrm{partners}(a)\}$
among its predecessors, so $r$ can only be enabled when every
named commit is done.

The same property is proved as a theorem in the companion Lean
library VegasCore as the ``Commit/reveal information barrier''
result over the canonical event graph
(\S\ref{sec:vegascore:theorems}): any event graph in which
commits dominate reveals satisfies the barrier, regardless of
how it was obtained.  Vegas's rewrite produces such graphs from
not-yet-canonical input; VegasCore certifies that any such graph
satisfies the barrier.

\subsection{Implementation}

The rewrite lives at \file{ir/EventGraph.kt} as a static method
\TT{expandCommitReveal} on the \TT{EventGraph} class.  It runs
once, immediately after the initial IR construction, before any
backend reads the graph.  The implementation is roughly sixty
lines of Kotlin and is reachability-driven: it iterates over the
existing nodes, builds the risk-partner relation by checking
$\neg(\mathrm{reach}(a,b) \vee \mathrm{reach}(b,a))$ for each
pair of pure-public nodes, and then performs the local rewrite
for each risk-partner-bearing node, allocating fresh node ids
above the current maximum.  The resulting graph is re-validated
by the EventGraph constructor (it must still be acyclic, the new
\COMMIT-\REVEAL pairs must be in order, and the read-visibility
discipline must be preserved).  The post-rewrite validation is
mostly defensive: all three properties are preserved by the
local rewrite construction, and the barrier property of
\S\ref{sec:cr:barrier} is preserved transitively.
